% !TeX root = ../main.tex

\appendix{A}{Prompt Templates}
\label{app:prompts}

This appendix reproduces the load-bearing prompt text of the system, so
that the mechanisms described in Chapters~\ref{ch:03} and~\ref{ch:04} can
be audited at the level at which they actually operate. Text is excerpted
verbatim from the repository (\texttt{agentic\_ip\_reuse/.../prompts.py}
and \texttt{ip\_reuse\_legacy/prompts.py}); ellipses mark elisions.

\section{Planner System Prompt (Catalog-Grounding Core)}

The planner's system prompt walks the model through the eight planning
stages and --- decisive in practice --- mandates the closed catalog
vocabulary:

\begin{lstlisting}[language={},caption={Planner system prompt, grounding
section (excerpt).},label={lst:planner-prompt}]
CATALOG GROUNDING (mandatory): The reusable-IP catalog for this
task is listed verbatim in the user message under "Reusable IP
catalog". Treat it as the only source of truth for which IPs
exist.
- Every reuse_decisions entry MUST set "selected_ip" to an
  "ip_id" that appears EXACTLY in that list.
- Never invent, rename, or add suffixes to an IP name (e.g. do
  not turn "e203_exu" into "e203_exu_core").
- If no catalog IP fits a module, set "new_rtl_required": true
  and "selected_ip": null for that module instead of guessing a
  name.
- Produce one reuse_decisions entry for every module that has a
  candidate IP in the catalog.
...
Return final output as one JSON object only, with keys:
requirements, modules, reuse_decisions, integration_plan,
verification_plan, debug_plan, unresolved_assumptions.
\end{lstlisting}

The user message then carries the injected catalog digest (one line per
provided module: \texttt{ip\_id}, port summary, one-sentence description)
and closes with: \emph{``Every reuse\_decisions.selected\_ip must be one
of the ip\_id values listed in the catalog above; do not invent
names.''} As Section~\ref{sec:results-grounding} showed, models comply
with the vocabulary yet never call the retrieval tools also offered ---
which is why the post-hoc grounding pass of
Section~\ref{sec:pipeline-planner} remains in place as the enforcement
backstop.

\section{Generation Prompt Structure}

The generator prompt is assembled from five blocks, in order: (1)~the
(possibly condensed) specification; (2)~the source of every reused module,
prefixed with the environment note that these modules \emph{are provided
as files and must be instantiated, never re-declared}; (3)~environment
notes (macros defined by the build, the one not-shipped dependency);
(4)~the testbench's DUT instantiation, verbatim, introduced as the binding
port contract; and (5)~the output-format instruction (one complete Verilog
module in a fenced code block).

\section{Functional Repair Prompt}

The functional repair turn (Section~\ref{sec:pipeline-funcrepair})
re-frames the task as logic repair on a design whose interface is already
correct. Its distinctive elements:

\begin{lstlisting}[language={},caption={Functional repair prompt, framing
(paraphrased structure; mismatch report inserted verbatim).},label={lst:func-prompt}]
The following {target_hdl} design COMPILES and matches the
required port interface, but its outputs MISMATCH the reference
testbench:

<verbatim simulator mismatch report, e.g.
 "Output 'b' has 23 mismatches. First mismatch at time 115.">

Fix the LOGIC only. Do not change the module name or the port
list. Reused modules listed as provided source files must only
be instantiated, never re-declared in your output.

<behavioral slice: specification sections mentioning the
 mismatching output identifiers, plus verbatim interface
 excerpts (Section 3.4.3); full spec if no identifier named>
\end{lstlisting}

When the repair loop runs on the direct flow (no plan), the reuse-framing
block is dropped and the same skeleton applies --- the mechanism by which
the router's direct arm shares the pipeline's repair machinery
(Section~\ref{sec:router-design}).
